\section{Teoremas de Isomorfismo}

\subsection*{Subgrupos}

\begin{definition}
    Un subconjunto no vacío \(S\subset G\) es \emph{subgrupo} \(S\leq G\) si verifica  
    \[\text{S}1. \ \ \forall s \in S, \ s^{-1} \in S. \hspace{2cm} \text{S}2. \ \ \forall s,t\in G, \ st\in G.\]
\end{definition}

\begin{theorem}
    \textbf{2.1.} \((S,*)\leq (G,*)\) es de facto grupo. \comentario{Via \(*|_{S \times S}\) + definición}
\end{theorem}

\begin{theorem}
    \textbf{2.2.} \(S\leq G \ \leftrightharpoons \ \id \in S \ \text{ e } \ \forall s,t \in S, \ st^{-1} \in S\).  
\end{theorem}

\begin{definition}
    Dado \(a \in G, \ \{a^m : m \in \Z\}=: \langle a \rangle \leq G\) es el \emph{subgrupo cíclico} generado por \(a\). En general, \(G\) es \emph{cíclico} si \(\exists a : G = \langle a \rangle\) e \(|\langle a \rangle|\) es el \emph{orden} de \(a\). \footnote{Usamos \(|\cdot |\) o \(\#\{\cdot\}\) para denotar cardinal, según quede más bonito.}
\end{definition}

\begin{theorem}
    \textbf{2.3.} Si \(|\langle a \rangle | = m < \infty\), entonces \(m = \min \{k : a^k = \id, \ k \in \N\}\). 
\end{theorem}

\demo{
    Es claro que \(\exists k \in \N : a^k = \id\). Sea \(\{a^j\}_{j=0}^{k-1} : \forall i < l \leq k-1, \ a^i \neq a^l\), ahí \(a^k = \id\) (por contradicción). Finalmente, \(\langle a \rangle = \{a^j\}\) (alg. de Euclides).  \hspace{-0.3cm}
}

\begin{corollary}
    \textbf{2.4.} Si \(|G| < \infty\), entonces \(S \leq G \ \leftrightharpoons \ \forall s,t \in S, \ st \in S\). \textcolor{gray}{(\(\Leftarrow\)) If \(s \in S\), then \(\langle s\rangle \leq S\), \(\exists m \in \N: s^m = \id\) and \(s^{-1} = s^{m-1} \in S\) }  
\end{corollary}

\begin{note}
    Claramente \(\id< S <G \leq G\), cuyo caso \(S\) es \emph{subgrupo propio}. 
\end{note}

\ejemplo{
    \centering Núcleo e Imagen de \(f: G \to H\) homomorfismo 
    \tcblower 
    \[\text{\scriptsize $\bullet $ \ \ }\{a \in G : f(a) = \id_{_H}\} =: \ker f\leq G;\hspace{5cm}\] 
    \[\hspace{4cm}\text{\scriptsize $\bullet $ \ \ }\{h \in H : h = f(a), \ a \in G\} =:\operatorname{Im}f\leq H.\]  
}

\hypertarget{thm25}{}
\begin{theorem}
    \textbf{2.5.} Sean \((S_i)_{i\in I} : S_i \leq G\), entonces \(\bigcap S_i \leq G\). \comentario{Mogollas }
\end{theorem}

%\begin{corollary}
%    \textbf{2.6.} If \(X \subset G, \ \exists H \leq G\), the smallest subgroup of \(G\) such that \( X \subset H\). \comentario{Let \(H := \)}
%\end{corollary}

\begin{definition}
    Dado \(X \subset G\), \(\langle X \rangle:= \bigcap \{S : S \leq G \text{ e } X \subset S\} \) es el subgrupo de \(G\) \emph{generado por} \(X\). En particular, si \(H, K \leq G\) denotamos \(H \lor K := \langle H \cup K \rangle  \). 
\end{definition}

\begin{note}
    Cuando \(|X|= n < \infty\), escribimos \(\langle X \rangle = \langle a_1, a_2, \ldots, a_n\rangle\).  
\end{note}

\begin{definition}
    Si \(\varnothing \neq X\subset G\), entonces una \emph{palabra} en \(X\) es una expresión de la forma \(x_1^{e_1} x_2^{e_2}\ldots x_n^{e_n}= w \in G\), donde \(x_j \in X\) y \(e_j = \pm 1 \).    
\end{definition}

\begin{theorem}
    \textbf{2.7.} \(\langle X \rangle \leq G\) es precisamente el conjunto de todas las palabras en \(X\). \comentario{Muestre que \(W:= \{\prod x_j^{e_j} : x_j \in X, \ e_j = \pm 1\}\leq G\) es el menor subgrupo \(: X \subset W\)}
\end{theorem}

\subsection*{Teorema de Lagrange}

\begin{definition}
    Sean \(S\leq G\) e \(t \in G\). La \emph{clase a derecha} de \(S\) en \(G\) es el subconjunto \(St := \{st : s \in S\}\), ahí \(t \) es un \emph{representante} de \(St\) (análogamente se define \(tS\)). 
\end{definition}

\ejemplo{
\begin{example}
    \(\langle n \rangle =: S < \Z\). Dado \(a \in \Z, \ a + S = \{a + qn : q \in \Z\} = [a] \in \Z_n \).   
\end{example}
}

\begin{note}
    En general \(tS\) e \(St\) no tienen porque coincidir. %, let \(H := \langle (1 \ \ 2) \rangle \leq S_3\). 
\end{note}

\hypertarget{lemma28}{}
\begin{lemma}
    \textbf{2.8.} \(Sa = Sb \ (\text{o }aS = bS) \ \leftrightharpoons \ ab^{-1} \in S \ (\text{o }b^{-1}a \in S)\). 
\end{lemma}

\demo{(\(\Rightarrow \)) \(\id a \in Sa = Sb\), luego \(\exists s \in S : a = sb \ \leftrightharpoons \ ab^{-1} = s \in S\). (\(\Leftarrow\)) Muestra las dos inclusiones usando que \(Sa \ni x = sa =  s\sigma b \in Sb: \sigma = a b^{-1}\).}

\hypertarget{thm29}{}
\begin{theorem}
    \textbf{2.9.} Cualesquier dos clases a derecha (o izquierda) de \(S\) en \(G\) son identicas o disyuntas. \comentario{Sigue del {\scriptsize\(\square\)} \hyperlink{lemma28}{2.8.} pues si \(x = sa = tb\), entonces \(ab^{-1} = s^{-1}t \in S\)}
\end{theorem}

\begin{theorem}
    \textbf{2.10.} El número de clases a derecha e izquierda \(S\) en \(G\) coincide. \comentario{Muestre que \(f: \mathscr{L} \to \mathscr{R}\) tal que \(Sa \mapsto a^{-1} S\) es una biyección bien definida}
\end{theorem}

\begin{definition}
    El \emph{índice} de \(S\) en \(G\), denotado por \([G:S]\), es el número de clases a derecha. 
\end{definition}

\hypertarget{thm211}{}
\begin{theorem}
    \textbf{2.11. (\emph{Lagrange})} Si \(|G| < \infty\) e \(S \leq G\), entonces \(|S| \) divide \( |G|\) e \([G:S] = \nicefrac{|G|}{|S|}\).  
\end{theorem}

\demo{
    Por \(\blacksquare\) \hyperlink{thm29}{2.9.} \(G =\bigsqcup^n St_j\), luego  \(|G| = \sum |St_j|\). Ahora, \(f_j : s\mapsto st_j\) es biyección bien definida, sigue que \(|St_j| = |S|\) e \(|G| = n |S|\), donde \(n = [G:S]\).   
}

\begin{corollary}
    \textbf{2.12.} Si \(|G|< \infty\) e \(a \in G\), entonces el orden de \(a\) divide \(|G|\). Más aún, si \(|G| = p \) primo, entonces \(|G|\) es cíclico. \comentario{Siendo \(a \neq \id \), \(|\langle a \rangle|\) divide \(p \ \leftrightharpoons \ |\langle a \rangle | = p\)}
\end{corollary}

\renewcommand{\mod}[1]{\ (\operatorname{mod} #1)}
\begin{corollary}
    \textbf{2.14. (\emph{Fermat})} Si \(p\) es primo e \(a \in \Z\), entonces \(a^p \equiv a \mod{p}\). 
\end{corollary}

\demo{Sean \(\F = \Z_p\) e \(G = \F^\times\). Si \([a] = [0]\), es claro que \([a]^p = [a] = [0] \). En otro caso, \([a]^{p-1} = [\id]\), luego \([a]^{p} = [a]\), de cualquier forma \(a^{p} \equiv a \mod{p}\).}

\subsection*{Grupos Cíclicos}

\hypertarget{lemma215}{}
\begin{lemma}
    \textbf{2.15.} Si \(G = \langle a\rangle : |G|= n < \infty\), entonces \(\exists ! S\leq G : |S| = d \) divide \(n\). 
\end{lemma}

\demo{Sea \(H = \langle b\rangle : |H|=d\) \ (\textcolor{rojoscuro}{Ex. 2.18}), entonces \(b^d = \id = {(a^{m})}^d\), luego \(\exists k \in \N : md = nk\) \ (\textcolor{rojoscuro}{Ex. 2.13}). Finalmente, \(b = a^m = (a^{\nicefrac{n}{d}})^k\) e \(H \leq \langle a^{\nicefrac{n}{d}}\rangle =: S \leq G\) (\textcolor{rojoscuro}{Ex. 2.11}). Como \( |S| = d \), tenemos que \(H= S\). }

\begin{definition}
    \centering \(\displaystyle \varphi(n) := \# \{k\in \N : k < n \text{ e } (k,n) = 1\}\).
\end{definition}
\hypertarget{thm216}{}
\begin{theorem}
    \textbf{2.16.} Si \(n\) es entero positivo, entonces \(\displaystyle n = \sum_{d \hspace{0.1cm} \mid \hspace{0.1cm}  n} \varphi(d)\).   
\end{theorem}

\demo{Sea \(G = \langle a\rangle : |G|=n\), entonces \(G = \bigsqcup \operatorname{gen}C : C \leq G\). Por {\scriptsize\(\square\)} \hyperlink{lemma215}{2.15.} \(\forall d \in \operatorname{Div}(n), \ \exists ! C_d \leq G\), luego \(|G| = \underset{d \ \mid \ n}{\sum} |\operatorname{gen}C_d| =\underset{d \ \mid \ n}{\sum}\varphi(d)\) (\textcolor{rojoscuro}{Ex. 2.20}). \vspace{-0.1cm} }

\begin{theorem}
    \textbf{2.17.} \(G : |G| = n\) es cíclico \(\leftrightharpoons \ \forall d : d  \mid  n \) se \(\exists  C_d \leq G : |C_d| = d\) es único.  
\end{theorem}

\demo{(\(\Rightarrow\)) {\scriptsize \(\square\) } \hyperlink{lemma215}{2.15.}. (\(\Leftarrow\)) Por el \(\blacksquare\) \hyperlink{thm216}{2.16.} \(n = \sum |\operatorname{gen}C| \leq \sum \varphi(d) = n\), entonces \(G\) debe tener subgrupo de orden \(n \ \leftrightharpoons \ G\) es cíclico. }

\begin{theorem}
    \textbf{2.18.} \(G\leq  \F^\times : |G| < \infty\) es cíclico. En particular, \(\F^\times\) es cíclico si \(|\F| < \infty \). \comentario{\(x^d - \id \) tiene a lo sumo \(d\) raízes en \(\F[x]\), el resultado sigue del \(\blacksquare\) \hyperlink{thm217}{2.17.} }
\end{theorem}

\begin{theorem}
    \textbf{2.19.} Sea \(p\) primo. Un grupo \(G : |G| = p^n\) es cíclico \ \(\leftrightharpoons\) \ es abeliano e \(\exists ! S\leq G : |S| = p\).  \comentario{Véase \cite[pag. 29]{stein}, requiere (\textcolor{rojoscuro}{Ex. 2.13}) }
\end{theorem}
%\demo{Let \(a \in G: |\langle a \rangle| = p^k\) is  largest as possible and \(S \leq \langle a \rangle \leq G\) the unique \(: |S| = p\). }

\subsection*{Subgrupos Normales}

\begin{definition}
    \centering \(ST := \{st : s \in S, \ t \in T\} \subset G\).   
\end{definition}

\begin{theorem}
    \textbf{2.20. (\emph{Fórmula del Producto})} Si \(S, T \leq G : |G|< \infty \), entonces \(|ST| |S \cap T| = |S||T|\). \footnote{\(ST\) no tiene porque ser subgrupo.} \comentario{No admite resumen, \cite[pag. 30]{stein}}
\end{theorem}

\begin{definition}
    \(K\leq G\) es \emph{subgrupo normal}, denotado por \(K\lhd G\), si \ \(\forall g \in G, \ gKg^{-1} = K\).  
\end{definition}

\ejemplo{
\begin{example}
    Si \(f : G \to H\) es homomorfismo, entonces \(\ker f \lhd G\). 
\end{example}
}

\begin{note}
    \(K \lhd G \ \leftrightharpoons \ Kg = gK\), luego existe una "conmutatividad parcial", si \(g \in G\) e \(k \in K, \ \exists k' \in K : kg = gk'\). 
\end{note}

\begin{definition}
    Un \emph{conjugado} de \(x\in G\) es un elemento de la forma \(axa^{-1} : a \in G\). 
\end{definition}

\begin{theorem}
    \textbf{2.21.} Si \(N \lhd G\), entonces el conjunto de clases \(\nicefrac{G}{N}\), es grupo (\emph{grupo cociente}) de orden \([G:N]\). \comentario{Haga \(Na * Nb := NaNb\) y demestre por definición }
\end{theorem}

\begin{corollary}
    \textbf{2.22.} Si \(N \lhd G\), entonces la \emph{proyección natural}, \(\pi: G \to \nicefrac{G}{N}\) tal que \( a\mapsto Na\), es homomorfismo sobrejectivo \(: \ker \pi = N\). \comentario{Mogollas}
\end{corollary}

\begin{note}
    El grupo cociente es la generalización de la construcción de \(\Z_n = \nicefrac{\Z}{\langle n \rangle}\) e todo \(N \lhd G\) es núcleo de algún homomorfismo.  
\end{note}

\begin{definition}
    El \emph{conmutador} de \(a,b \in G\) es \([a,b]:= aba^{-1}b^{-1}\) y el \emph{subgrupo conmutador} (\emph{derivado}) es \(G':= \langle X\rangle\), donde \(X :=\{ [a,b] : a,b \in G\}\). 
\end{definition}

\alerta{\begin{example}
    \(X := \{[a,b] : a,b \in G\}\) no necesariamente es subgrupo, \([a,b] * [c,d]\) no tiene porque ser conmutador (\textcolor{rojoscuro}{Ex. 2.43}).
\end{example} 
    }

\hypertarget{thm223}{}
\begin{theorem}
    \textbf{2.23.} \(G' \lhd G\). Más aún, si \(H\lhd G\), entonces \(\nicefrac{G}{H}\) es abeliano \ \(\leftrightharpoons\) \ \(G' \leq H\). 
\end{theorem}

\demo{Si \(f : G \to G \) es homomorfismo, entonces \(f([a,b]) = [f(a), f(b)]\) luego \(G' \lhd G\) (\textcolor{rojoscuro}{Ex. 2.34}). Más aún, \(Hab = Hba \ \leftrightharpoons \ (ab)(ba)^{-1} = a b a^{-1}b^{-1} = [a,b] \in H\).}

\subsection*{Los Teoremas de Isomorfismo}

%\parbox[a]{0.8\linewidth}{
\hypertarget{thm224}{}
\begin{theorem}
    \textbf{2.24. (\emph{\(1^\text{er}\) Teorema de Isomorfismo})} Sea \(f : G \to H\) homomorfismo con \(\ker f = K \). Entonces, \(K \lhd G\) e \(\exists \varphi : \nicefrac{G}{K} \cong \operatorname{Im}f\). \comentario{Haga \(\varphi : Ka \mapsto f(a)\) }
\end{theorem}
%}

\begin{lemma}
    \textbf{2.25.} Si \(S, T \leq G\) e \(S \lhd G\) (o \(T\lhd G\)), entonces \(ST = S \lor T = TS\). 
\end{lemma}

\demo{
    (\(\subseteq \)) Mogollas; (\(\supseteq\)) Sigue si \(ST, TS  \leq G\), de la nota después del \(\blacksquare\) \hyperlink{thm25}{2.5.}. Si \(T \lhd G, \ s_j\in S\) e \(t_j \in T\) muestre que \(s_1t_1(s_2t_2)^{-1} \in ST\).  
}

\hypertarget{thm226}{}
\begin{theorem}
    \textbf{2.26. (\emph{\(2^\text{do}\) Teorema de Isomorfismo})} Sean \(T, N \leq G\) con \(N\) normal. Entonces, \(N \cap T\lhd T\) e \(\nicefrac{T}{N\cap T}\cong \nicefrac{NT}{N}\). \comentario{Sigue del \(\blacksquare\) \hyperlink{thm224}{2.24.} con \(\ker \pi|_T = N \cap T \lhd T\) e \(\nicefrac{T}{N\cap T} \cong \operatorname{Im}\pi|_T = \nicefrac{NT}{N}\)}
\end{theorem}

\begin{theorem}
    \textbf{2.27. (\emph{\(3^\text{er}\) Teorema de Isomorfismo})} Sean \(K \leq H\) ambos \(K, H \lhd G\). Entonces, \(\nicefrac{H}{K}\lhd \nicefrac{G}{K}\) e \((\nicefrac{G}{K}) \ \slash \ (\nicefrac{H}{K}) \cong \nicefrac{G}{H}\). \comentario{Sigue del \(\blacksquare\) \hyperlink{thm224}{2.24.}. con \(f: \nicefrac{G}{K} \twoheadrightarrow \nicefrac{G}{H}\  \mid \  Ka \mapsto Ha\) e \(\ker f = \nicefrac{H}{K}\) }
\end{theorem}

\hypertarget{thm228}{}
\begin{theorem}
    \textbf{2.28. (\emph{Correspondencia})} Sean \(K \lhd G\) e \(\pi : G \to \nicefrac{G}{K}\). Entonces, \(S \mapsto \pi(S) = \nicefrac{S}{K}\) es una biyección entre \(\{S: K \leq S \leq G\}\) e \(\{H : H \leq \nicefrac{G}{K}\}\). Más aún, si \(S^* = \nicefrac{S}{K}\), 
    \begin{enumerate}[label = \roman*., left = 0.35cm]
        \item \vspace{-0.2cm}\(T \leq S \ \leftrightharpoons \ T^* \leq S^*\), luego \([S:T] = [S^*: T^*]\). 
        \item \(T \lhd S \ \leftrightharpoons \ T^* \lhd S^*\), luego \(\nicefrac{S}{T}= \nicefrac{S^*}{T^*}\). 
    \end{enumerate}
\end{theorem}

\vspace{-0.2cm}
\comentario{No soporta resumen, véase \cite[pag. 38]{stein}}

\begin{definition}
    Un grupo \(\id < G\) es \emph{simple} si no contiene subgrupos propios normales.  
\end{definition}

\subsection*{Productos Directos}

\begin{definition}
    Dados \(H\) e \( K \) grupos, su \emph{producto directo}, denotado por \(H \times K\) es un grupo con la operación (por coordenadas), \ \(*: (h,k) \times (h',k') \mapsto (h h', k k')\). 
\end{definition}

\begin{note}
    \(H \cong H \times \id_{_K}\) and \(K \cong \id_{_H} \times K\). 
\end{note}

\hypertarget{thm229}{}
\begin{theorem}
    \textbf{2.29.} Sean \(H, K \lhd G\). Si \(HK = G \) e \(H \cap K = \id\), entonces \(G \cong H \times K\).  
\end{theorem}

\demo{Defina \(f: G \to H \times K  \) tal que \( a = hk \mapsto (h,k)\), use conmutadores para ver que \(h'kh'^{-1}k^{-1} \in H \cap K\) y que es homomorfismo biyectivo bien definido. }

\alerta{ \centering Todas las hipótesis en el \(\blacksquare\) \hyperlink{thm229}{2.29.} son necesarias
\tcblower
\begin{example}
    \(G = S_3\) (no abeliano), \(H = \langle (1 \ \ 2 \ \ 3)\rangle \) e \(K = \langle (1 \ \ 2) \rangle\). Verifican todo salvo que \(K\ntriangleleft G\). Tenemos \(H \times K \cong \Z_3 \times \Z_2\) (abeliano).\Lightning 
\end{example}
}

\begin{theorem}
    \textbf{2.30.} Si \(A \lhd H\) e \(B \lhd K\), entonces \(A \times B \lhd H \times K\) e \ \(\nicefrac{H\times K}{A \times B} \cong \nicefrac{H}{A}\times \nicefrac{K}{B}\). 
\end{theorem}

\demo{Sigue del \(\blacksquare\) \hyperlink{thm224}{2.24.} con \(\varphi : H\times K \to \nicefrac{H}{A} \times \nicefrac{K}{B}\) (sobreyectiva) tal que \((h,k) \mapsto (Ah, Bk)\), donde \(\ker \varphi = A \times B\). 
} 

\begin{corollary}
    \textbf{2.31.} Si \(G = H \times K \), entonces \(\nicefrac{G}{H\times \id}\cong \nicefrac{H}{H} \times \nicefrac{K}{\id} \cong K\).  
\end{corollary}